\documentclass[12pt,a4paper]{article}

% ── Encoding & Language ──
\usepackage[utf8]{inputenc}
\usepackage[spanish,es-tabla]{babel}

% ── Fonts ──
\usepackage{libertine}           % Linux Libertine — elegante para cuerpo
\usepackage{tgheros}              % TeX Gyre Heros (Helvetica) — headers
\usepackage{sourcesanspro}        % Source Sans Pro — tablas y captions
\usepackage[T1]{fontenc}

% ── Layout ──
\usepackage{geometry}
\geometry{margin=2.2cm, headheight=18pt, top=2.5cm, bottom=2.2cm}

% ── Colors — paleta ingenieril ──
\usepackage{xcolor}
\definecolor{primary}{HTML}{1B3A5C}    % azul oscuro ingenieril
\definecolor{secondary}{HTML}{2E86AB}  % azul medio
\definecolor{accent}{HTML}{E85D2C}     % naranja calido
\definecolor{lightbg}{HTML}{EDF2F7}    % gris azulado claro
\definecolor{tablerow}{HTML}{F0F4F8}   % alternado tablas
\definecolor{darktext}{HTML}{1A202C}   % texto oscuro
\definecolor{softrule}{HTML}{CBD5E0}   % lineas suaves

% ── Headers/Footers ──
\usepackage{fancyhdr}
\pagestyle{fancy}
\fancyhf{}
\fancyhead[L]{\sffamily\footnotesize\color{primary!80}MEMORIA DE CALCULO ESTRUCTURAL}
\fancyhead[R]{\sffamily\footnotesize\color{primary!80}Cercha Warren Cajon 3D}
\fancyfoot[C]{\sffamily\footnotesize\color{primary!60}\thepage}
\renewcommand{\headrulewidth}{0.4pt}
\renewcommand{\headrule}{\hbox to\headwidth{\color{secondary!40}\leaders\hrule height \headrulewidth\hfill}}

% ── Section styling ──
\usepackage{titlesec}
\titleformat{\section}[hang]
  {\sffamily\LARGE\bfseries\color{primary}}
  {\thesection\enskip}
  {0pt}
  {}
  [\vspace{-0.5em}\color{accent}{\rule{\textwidth}{1.2pt}}\vspace{0.2em}]
\titleformat{\subsection}[hang]
  {\sffamily\Large\bfseries\color{secondary}}
  {\thesubsection\enskip}
  {0pt}
  {}

% ── Tables ──
\usepackage{booktabs}
\usepackage{longtable}
\usepackage{colortbl}
\usepackage{array}
\newcommand{\tableheadrow}{\rowcolor{primary!15}\sffamily\footnotesize\bfseries\color{primary}}
\newcommand{\tabletotalrow}{\rowcolor{secondary!10}\sffamily\bfseries}

% ── Figures & Graphics ──
\usepackage{graphicx}
\usepackage{tikz}
\usetikzlibrary{babel}
\usepackage{float}

% ── Verbatim for appendix ──
\usepackage{verbatim}

% ── Hyperlinks ──
\usepackage{hyperref}
\hypersetup{
  colorlinks=true,
  linkcolor=secondary,
  urlcolor=accent,
  citecolor=secondary,
  pdfborderstyle={/S/U/W 0.5},
}

% ── Captions ──
\usepackage{caption}
\captionsetup{
  font={sf,footnotesize,color=primary!80},
  labelfont={bf,color=secondary},
  justification=raggedright,
  singlelinecheck=false,
  skip=6pt,
}

% ── Metadata ──
\hypersetup{
  pdftitle={Memoria de Calculo — Cercha Warren Cajon 3D},
  pdfauthor={DogCalC v0.4.2 — Ing. Civil},
  pdfsubject={Analisis Estructural NSR-10},
  pdfkeywords={cercha, warren, acero, NSR-10, AISC 360-10, LRFD},
}

\begin{document}

% ═══════════════════════════════════════════
%  TITLE PAGE
% ═══════════════════════════════════════════
\thispagestyle{empty}
\begin{tikzpicture}[remember picture,overlay]
  % Barra superior
  \fill[primary] (current page.north west) rectangle ([yshift=-2.8cm]current page.north east);
  % Barra inferior
  \fill[secondary!80] (current page.south west) rectangle ([yshift=1.2cm]current page.south east);
  % Linea acento
  \fill[accent] (current page.north west) rectangle ([yshift=-2.85cm]current page.north east);
\end{tikzpicture}

\vspace*{2cm}
\begin{center}
  {\sffamily\fontsize{16}{20}\selectfont\color{white}\bfseries
   MEMORIA DE CALCULO ESTRUCTURAL\par}
  \vspace{0.8cm}
  {\sffamily\fontsize{28}{34}\selectfont\color{primary}\bfseries
   Cercha Warren\par}
  {\sffamily\fontsize{22}{28}\selectfont\color{secondary}\bfseries
   Cajon 3D — Cubierta Metalica\par}
  \vspace{1.2cm}
  {\sffamily\large\color{darktext!70}
   Analisis por Elementos Finitos\par}
  \vspace{0.3cm}
  {\sffamily\large\color{darktext!70}
   Norma NSR-10 — Diseno LRFD (AISC 360-10)\par}
  \vspace{2.5cm}

  \begin{tikzpicture}[scale=1.2]
    % Cercha esquematica
    \draw[primary!50, thick] (0,0.8) -- (8,0.8);
    \draw[primary!50, thick] (0,0) -- (8,0);
    \foreach \x in {0,1,2,3,4,5,6,7,8} {
      \draw[secondary!30, thin] (\x,0) -- (\x,0.8);
    }
    \foreach \x in {0,1,...,7} {
      \draw[accent!40, thin] (\x,0.8) -- (\x+1,0);
    }
    \node[primary!40] at (4,1.05) {\sffamily\scriptsize 16.52 m};
    \node[secondary!40] at (8.3,0.4) {\sffamily\scriptsize 1.0 m};
  \end{tikzpicture}

  \vspace{2cm}
  {\sffamily\large\color{white}\bfseries DogCalC v0.4.2 — Motor PyNite FEM\par}
  \vspace{0.3cm}
  {\sffamily\small\color{white!80}\today\par}
\end{center}

\newpage

% ═══════════════════════════════════════════
%  TABLE OF CONTENTS
% ═══════════════════════════════════════════
\pagestyle{fancy}
\renewcommand{\contentsname}{\sffamily\color{primary}\LARGE\bfseries Contenido}
\tableofcontents
\newpage

%============================================
\section{DATOS GENERALES}
%============================================

\begin{center}
\begin{tabular}{>{\sffamily\color{primary!80}}l >{\sffamily}l}
\textbf{Proyecto:}   & Cercha Warren en cajon 3D para cubierta \\
\textbf{Software:}   & DogCalC v0.4.2 (motor PyNite FEM) \\
\textbf{Norma:}      & NSR-10 (AISC 360-10 para acero) \\
\textbf{Metodo:}     & Diseno por Resistencia Ultima (LRFD) \\
\end{tabular}
\end{center}

%============================================
\section{GEOMETRIA}
%============================================

\subsection{Vista lateral -- Cara A (Z=0)}

\begin{figure}[H]
\centering
\begin{tikzpicture}[scale=0.55, every node/.style={font=\tiny\sffamily}]
  \draw[thick, primary] (0,1.0) -- (16.52,1.0);
  \draw[thick, accent] (0,0) -- (16.52,0);
  \draw[thin, softrule] (0.00,0) -- (0.00,1.0);
  \fill[primary] (0.00,1.0) circle (0.06) node[above]{\tiny 1};
  \fill[accent] (0.00,0) circle (0.06) node[below]{\tiny 21};
  \draw[thin, softrule] (0.90,0) -- (0.90,1.0);
  \fill[primary] (0.90,1.0) circle (0.06) node[above]{\tiny 2};
  \fill[accent] (0.90,0) circle (0.06) node[below]{\tiny 22};
  \draw[thin, softrule] (1.80,0) -- (1.80,1.0);
  \fill[primary] (1.80,1.0) circle (0.06) node[above]{\tiny 3};
  \fill[accent] (1.80,0) circle (0.06) node[below]{\tiny 23};
  \draw[thin, softrule] (2.70,0) -- (2.70,1.0);
  \fill[primary] (2.70,1.0) circle (0.06) node[above]{\tiny 4};
  \fill[accent] (2.70,0) circle (0.06) node[below]{\tiny 24};
  \draw[thin, softrule] (3.60,0) -- (3.60,1.0);
  \fill[primary] (3.60,1.0) circle (0.06) node[above]{\tiny 5};
  \fill[accent] (3.60,0) circle (0.06) node[below]{\tiny 25};
  \draw[thin, softrule] (4.50,0) -- (4.50,1.0);
  \fill[primary] (4.50,1.0) circle (0.06) node[above]{\tiny 6};
  \fill[accent] (4.50,0) circle (0.06) node[below]{\tiny 26};
  \draw[thin, softrule] (5.40,0) -- (5.40,1.0);
  \fill[primary] (5.40,1.0) circle (0.06) node[above]{\tiny 7};
  \fill[accent] (5.40,0) circle (0.06) node[below]{\tiny 27};
  \draw[thin, softrule] (6.30,0) -- (6.30,1.0);
  \fill[primary] (6.30,1.0) circle (0.06) node[above]{\tiny 8};
  \fill[accent] (6.30,0) circle (0.06) node[below]{\tiny 28};
  \draw[thin, softrule] (7.13,0) -- (7.13,1.0);
  \fill[primary] (7.13,1.0) circle (0.06) node[above]{\tiny 9};
  \fill[accent] (7.13,0) circle (0.06) node[below]{\tiny 29};
  \draw[thin, softrule] (8.03,0) -- (8.03,1.0);
  \fill[primary] (8.03,1.0) circle (0.06) node[above]{\tiny 10};
  \fill[accent] (8.03,0) circle (0.06) node[below]{\tiny 30};
  \draw[thin, softrule] (8.93,0) -- (8.93,1.0);
  \fill[primary] (8.93,1.0) circle (0.06) node[above]{\tiny 11};
  \fill[accent] (8.93,0) circle (0.06) node[below]{\tiny 31};
  \draw[thin, softrule] (9.83,0) -- (9.83,1.0);
  \fill[primary] (9.83,1.0) circle (0.06) node[above]{\tiny 12};
  \fill[accent] (9.83,0) circle (0.06) node[below]{\tiny 32};
  \draw[thin, softrule] (10.73,0) -- (10.73,1.0);
  \fill[primary] (10.73,1.0) circle (0.06) node[above]{\tiny 13};
  \fill[accent] (10.73,0) circle (0.06) node[below]{\tiny 33};
  \draw[thin, softrule] (11.55,0) -- (11.55,1.0);
  \fill[primary] (11.55,1.0) circle (0.06) node[above]{\tiny 14};
  \fill[accent] (11.55,0) circle (0.06) node[below]{\tiny 34};
  \draw[thin, softrule] (12.45,0) -- (12.45,1.0);
  \fill[primary] (12.45,1.0) circle (0.06) node[above]{\tiny 15};
  \fill[accent] (12.45,0) circle (0.06) node[below]{\tiny 35};
  \draw[thin, softrule] (13.35,0) -- (13.35,1.0);
  \fill[primary] (13.35,1.0) circle (0.06) node[above]{\tiny 16};
  \fill[accent] (13.35,0) circle (0.06) node[below]{\tiny 36};
  \draw[thin, softrule] (14.25,0) -- (14.25,1.0);
  \fill[primary] (14.25,1.0) circle (0.06) node[above]{\tiny 17};
  \fill[accent] (14.25,0) circle (0.06) node[below]{\tiny 37};
  \draw[thin, softrule] (15.15,0) -- (15.15,1.0);
  \fill[primary] (15.15,1.0) circle (0.06) node[above]{\tiny 18};
  \fill[accent] (15.15,0) circle (0.06) node[below]{\tiny 38};
  \draw[thin, softrule] (16.05,0) -- (16.05,1.0);
  \fill[primary] (16.05,1.0) circle (0.06) node[above]{\tiny 19};
  \fill[accent] (16.05,0) circle (0.06) node[below]{\tiny 39};
  \draw[thin, softrule] (16.52,0) -- (16.52,1.0);
  \fill[primary] (16.52,1.0) circle (0.06) node[above]{\tiny 20};
  \fill[accent] (16.52,0) circle (0.06) node[below]{\tiny 40};

  \foreach \x/\xnext in {0.00/0.90, 0.90/1.80, 1.80/2.70, 2.70/3.60, 3.60/4.50, 4.50/5.40, 5.40/6.30, 6.30/7.13, 7.13/8.03, 8.03/8.93, 8.93/9.83, 9.83/10.73, 10.73/11.55, 11.55/12.45, 12.45/13.35, 13.35/14.25, 14.25/15.15, 15.15/16.05, 16.05/16.52} {
    \draw[thin, secondary!40] (\x,1.0) -- (\xnext,0);
  }

  \node[primary, font=\tiny\sffamily] at (0.45,1.18) {M1 (CS)};
  \node[accent, font=\tiny\sffamily] at (0.45,{-0.28}) {M20 (CI)};
  \draw[<->, softrule] (0,-0.5) -- (16.52,-0.5) node[midway,below,font=\tiny\sffamily] {Luz: 16.52 m};
  \draw[<->, softrule] (16.52+0.5,0) -- (16.52+0.5,1.0) node[midway,right,font=\tiny\sffamily] {h=1.0m};
  \node[primary, right, font=\tiny\sffamily] at (0,-0.9) {CS = Cordon Superior (compresion)};
  \node[accent, right, font=\tiny\sffamily] at (7,-0.9) {CI = Cordon Inferior (traccion)};
\end{tikzpicture}
\caption{Vista lateral de la cercha Warren. Nodos 1--20 en cordon superior,
nodos 21--40 en cordon inferior. Diagonales en celeste, montantes en gris.}
\label{fig:vista-lateral}
\end{figure}

\subsection{Seccion transversal -- Cajon 3D}

\begin{figure}[H]
\centering
\begin{tikzpicture}[scale=1.5, every node/.style={font=\tiny\sffamily}]
  \draw[thick, primary] (0,0) rectangle (0.3,1.0);
  \draw[thick, primary] (2.0,0) rectangle (2.3,1.0);
  \draw[thin, dashed, softrule] (0,0) -- (2.0,0);
  \draw[thin, dashed, softrule] (0,1.0) -- (2.0,1.0);
  \draw[thin, dashed, softrule] (0.3,0) -- (2.3,0);
  \draw[thin, dashed, softrule] (0.3,1.0) -- (2.3,1.0);
  \fill[primary] (0,0) circle (0.05) node[left] {\tiny N21};
  \fill[primary] (0,1.0) circle (0.05) node[left] {\tiny N1};
  \fill[primary] (0.3,0) circle (0.05) node[right] {\tiny N22};
  \fill[primary] (0.3,1.0) circle (0.05) node[right] {\tiny N2};
  \fill[primary] (2.0,0) circle (0.05) node[left] {\tiny N61};
  \fill[primary] (2.0,1.0) circle (0.05) node[left] {\tiny N41};
  \fill[primary] (2.3,0) circle (0.05) node[right] {\tiny N62};
  \fill[primary] (2.3,1.0) circle (0.05) node[right] {\tiny N42};
  \draw[<->, accent] (0.3,0.5) -- (2.0,0.5) node[midway,above] {\tiny 0.28 m};
  \node at (0.15,-0.3) {\tiny Cara A (Z=0)};
  \node at (2.15,-0.3) {\tiny Cara B (Z=0.28)};

  \node[right] at (3.0,1.0) {\sffamily\bfseries\color{primary} TUBO 70$\times$70$\times$4};
  \node[right] at (3.0,0.7) {\sffamily\footnotesize $A = 1{,}056$ mm$^2$};
  \node[right] at (3.0,0.4) {\sffamily\footnotesize $I = 769{,}472$ mm$^4$};
  \node[right] at (3.0,0.1) {\sffamily\footnotesize $r = 27.0$ mm};
\end{tikzpicture}
\caption{Seccion transversal del cajon. Dos cerchas paralelas separadas 0.28 m,
unidas por montantes transversales en cada nudo.}
\label{fig:seccion}
\end{figure}

\subsection{Datos geometricos}

\begin{center}
\begin{tabular}{>{\sffamily\color{primary!80}}l >{\sffamily}l}
\textbf{Tipo:}              & Cercha Warren en cajon 3D (dos caras paralelas) \\
\textbf{Luz total:}         & 16.52 m \\
\textbf{Altura:}            & 1.00 m \\
\textbf{Separacion caras:} & 0.28 m (ancho del cajon) \\
\textbf{Numero de paneles:} & 19 \\
\textbf{Nodos totales:}     & 80 (40 por cara) \\
\textbf{Miembros totales:}  & 194 \\
\end{tabular}
\end{center}

%============================================
\section{MATERIALES Y SECCIONES}
%============================================

{
\sffamily
\begin{tabular}{>{\color{primary!80}}l l}
\textbf{Acero:}        & ASTM A36 equivalente \\
\textbf{Modulo E:}     & $E = 200\,000$ MPa \\
\textbf{Poisson:}      & $\nu = 0.30$ \\
\textbf{Densidad:}     & $\rho = 7\,850$ kg/m\textsuperscript{3} \\
\textbf{Fluencia:}     & $F_y = 250$ MPa \\
\textbf{Ultimo:}       & $F_u = 400$ MPa \\
\end{tabular}
}

\subsection{Perfil}

\begin{center}
\begin{tabular}{l>{\sffamily}c}
\toprule
\rowcolor{primary!15}
\multicolumn{1}{c}{\sffamily\footnotesize\bfseries\color{primary} Propiedad} & \multicolumn{1}{c}{\sffamily\footnotesize\bfseries\color{primary} Valor} \\
\midrule
Designacion   & TUBO 70$\times$70$\times$4 mm \\
Area ($A_g$)   & 1\,056 mm\textsuperscript{2} \\
Inercia ($I$)   & 769\,472 mm\textsuperscript{4} \\
Radio de giro ($r$) & 27.0 mm \\
\bottomrule
\end{tabular}
\end{center}

%============================================
\section{CARGAS}
%============================================

{
\sffamily
Cargas aplicadas en nudos del cordon superior de ambas caras,
ancho aferente de \textbf{1.30 m} (0.65 m a cada lado).
}

\vspace{0.4cm}
\begin{center}
\begin{tabular}{l c >{\sffamily}c}
\toprule
\tableheadrow
\textbf{Carga} & \textbf{kN/m\textsuperscript{2}} & \textbf{Total (kN)} \\
\midrule
Peso propio (D)     & 0.80  & 34.36 \\
Cubierta viva (Lr)  & 0.50  & 21.48 \\
\rowcolor{tablerow}
Viento (W)          & 0.40  & 17.18 \\
Granizo (G)         & 1.00  & 42.95 \\
\bottomrule
\end{tabular}
\end{center}

\subsection{Casos de carga}

{
\sffamily
\begin{enumerate}
    \item \textbf{DEAD} -- Peso propio + carga muerta (D)
    \item \textbf{ROOF LIVE} -- Carga viva de cubierta (Lr)
    \item \textbf{WIND} -- Viento (W)
    \item \textbf{GRANIZO} -- Granizo (G)
\end{enumerate}

Todos incluyen peso propio estructural (SELFWEIGHT Y -1).
}

%============================================
\section{COMBINACIONES NSR-10}
%============================================

Combinaciones de Resistencia Ultima (sin sismo):

\vspace{0.3cm}
\begin{center}
\begin{tabular}{c >{\sffamily}l >{\sffamily}l}
\toprule
\tableheadrow
\textbf{Combo} & \textbf{Nombre} & \textbf{Factores} \\
\midrule
101 & 1.4D          & $1.4\,D$ \\
\rowcolor{tablerow}
102 & 1.2D+1.6Lr     & $1.2\,D + 1.6\,Lr$ \\
103 & 1.2D+1.6G      & $1.2\,D + 1.6\,G$ \\
\rowcolor{tablerow}
104 & 1.2D+1.6Lr+0.5W & $1.2\,D + 1.6\,Lr + 0.5\,W$ \\
105 & 1.2D+1.6G+0.5W  & $1.2\,D + 1.6\,G + 0.5\,W$ \\
\rowcolor{tablerow}
106 & 1.2D+1.0W+0.5Lr & $1.2\,D + 1.0\,W + 0.5\,Lr$ \\
107 & 1.2D+1.0W+0.5G  & $1.2\,D + 1.0\,W + 0.5\,G$ \\
\rowcolor{tablerow}
108 & 0.9D+1.0W       & $0.9\,D + 1.0\,W$ \\
\bottomrule
\end{tabular}
\end{center}

Envolvente de diseno: ENVELOPE 1 TYPE STRENGTH (combos 101--108).

%============================================
\section{ANALISIS ESTRUCTURAL}
%============================================

{
\sffamily
Modelo 3D: 80 nudos, 194 elementos frame (6 GDL/nudo), conexiones rigidas.

\textbf{Apoyos:}
\begin{itemize}
    \item N21 y N61 (izquierda): Fijo (PINNED)
    \item N40 y N80 (derecha): Deslizante en X (ROLLER)
\end{itemize}
}

\subsection{Reacciones -- Combo 105 (1.2D+1.6G+0.5W)}

\begin{center}
\begin{tabular}{l r >{\sffamily}r r}
\toprule
\tableheadrow
\textbf{Nudo} & \textbf{FX (kN)} & \textbf{FY (kN)} & \textbf{FZ (kN)} \\
\midrule
N21 (Fijo)     & --    & 36.7  & -- \\
\rowcolor{tablerow}
N61 (Fijo)     & --    & 36.7  & -- \\
N40 (Rodillo)  & 0     & 36.3  & -- \\
\rowcolor{tablerow}
N80 (Rodillo)  & 0     & 36.3  & -- \\
\midrule
\tabletotalrow
\textbf{Total} & 0     & \textbf{145.9} & 0 \\
\bottomrule
\end{tabular}
\end{center}

{
\sffamily
Desplazamiento maximo: $\delta_{max} \approx$ \textbf{42.9 mm} $\approx L / 385$.
}

%============================================
\section{FUERZAS AXIALES}
%============================================

Envolvente de diseno, Cara A (Z=0). $(+)$ traccion, $(-)$ compresion. Valores en kN.

\subsection{Cordon Superior (compresion)}

\begin{center}
\begin{tabular}{r >{\sffamily}r | r >{\sffamily}r}
\toprule
\tableheadrow
\textbf{Miembro} & \textbf{Axial (kN)} & \textbf{Miembro} & \textbf{Axial (kN)} \\
\midrule
  M1  &   -30.4 & M11 &  -147.9 \\
\rowcolor{tablerow}
  M2  &   -57.7 & M12 &  -143.6 \\
  M3  &   -81.5 & M13 &  -135.7 \\
\rowcolor{tablerow}
  M4  &  -101.6 & M14 &  -125.3 \\
  M5  &  -118.2 & M15 &  -110.6 \\
\rowcolor{tablerow}
  M6  &  -131.3 & M16 &   -92.3 \\
  M7  &  -140.8 & M17 &   -70.5 \\
\rowcolor{tablerow}
  M8  &  -146.3 & M18 &   -45.2 \\
  M9  &  -149.0 & M19 &   -15.8 \\
\rowcolor{tablerow}
  M10 &  -148.1 \\
\bottomrule
\end{tabular}
\end{center}

{
\sffamily
\textbf{Maxima compresion:} M9 = \textcolor{accent}{-149.0 kN} (centro de la luz).
}

\subsection{Cordon Inferior (traccion)}

\begin{center}
\begin{tabular}{r >{\sffamily}r | r >{\sffamily}r}
\toprule
\tableheadrow
\textbf{Miembro} & \textbf{Axial (kN)} & \textbf{Miembro} & \textbf{Axial (kN)} \\
\midrule
  M20  &    -4.9 & M30 &   143.1 \\
\rowcolor{tablerow}
  M21  &    26.2 & M31 &   135.4 \\
  M22  &    53.9 & M32 &   125.1 \\
\rowcolor{tablerow}
  M23  &    78.0 & M33 &   110.6 \\
  M24  &    98.6 & M34 &    92.5 \\
\rowcolor{tablerow}
  M25  &   115.6 & M35 &    70.8 \\
  M26  &   129.0 & M36 &    45.6 \\
\rowcolor{tablerow}
  M27  &   138.9 & M37 &    16.6 \\
  M28  &   144.8 & M38 &     0.7 \\
\rowcolor{tablerow}
  M29  &   147.9 \\
\bottomrule
\end{tabular}
\end{center}

{
\sffamily
\textbf{Maxima traccion:} M29 = \textcolor{secondary}{+147.9 kN} (centro de la luz).
}

\subsection{Diagonales y Montantes}

\begin{center}
\begin{tabular}{r >{\sffamily}r r | r >{\sffamily}r r}
\toprule
\rowcolor{primary!15}
\multicolumn{3}{c}{\sffamily\footnotesize\bfseries\color{primary} Diagonales} & \multicolumn{3}{c}{\sffamily\footnotesize\bfseries\color{primary} Montantes} \\
\midrule
M39  & 44.8  & (T) & M58  & --35.8 & (C) \\
\rowcolor{tablerow}
M116 & 44.8  & (T) & M135 & --35.8 & (C) \\
M56  & 42.2  & (T) & M77  & --34.1 & (C) \\
\rowcolor{tablerow}
M133 & 42.2  & (T) & M154 & --34.1 & (C) \\
M40  & 40.0  & (T) & M59  & --33.1 & (C) \\
\rowcolor{tablerow}
M117 & 40.0  & (T) & M136 & --33.1 & (C) \\
M55  & 36.7  & (T) & M76  & --31.9 & (C) \\
\bottomrule
\end{tabular}
\end{center}

%============================================
\section{CODE CHECK NSR-10}
%============================================

\subsection{Capacidades}

\begin{center}
\begin{tabular}{l >{\sffamily}l}
\toprule
\tableheadrow
\textbf{Modo de falla} & \textbf{$\phi P_n$} \\
\midrule
Traccion (AISC D2):   & $\phi_t P_n = 0.9 \times 250 \times 1056 = \mathbf{238}$ kN \\
Compresion (AISC E3): & $\phi_c P_n = 0.9 \times F_{cr} \times A_g = \mathbf{224}$ kN \\
\bottomrule
\end{tabular}
\end{center}

\subsection{Resultados}

\begin{center}
\begin{tabular}{l >{\sffamily}r}
\toprule
\tableheadrow
\textbf{Concepto} & \textbf{Valor} \\
\midrule
Total miembros verificados  & 194 \\
Miembros OK                 & \textcolor{secondary!80!black}{194} \\
\rowcolor{tablerow}
Miembros SOBREESFORZADOS    & \textcolor{accent}{\textbf{0}} \\
Peor ratio demanda/capacidad & \textbf{0.67} (M9, cordon superior) \\
\bottomrule
\end{tabular}
\end{center}

\vspace{0.5cm}
{
\sffamily\small
\color{primary!80}
\begin{center}
\begin{minipage}{0.85\textwidth}
\rule{\textwidth}{0.3pt}\\[4pt]
\textbf{Conclusion:} Con el ancho aferente reducido a 0.65 m, el perfil
TUBO 70$\times$70$\times$4 mm \textbf{pasa holgadamente} el code check
(ratio maximo = 0.67). Los 194 miembros cumplen AISC 360-10 sin sobreesfuerzos.\\[2pt]
\rule{\textwidth}{0.3pt}
\end{minipage}
\end{center}
}

%============================================
\section{CONCLUSIONES}
%============================================

{
\sffamily
\begin{enumerate}
    \item La cercha Warren en cajon 3D es \textbf{estructuralmente estable}.
          El analisis FEM confirma equilibrio para las 8 combinaciones NSR-10.

    \item Con el ancho aferente de 0.65 m, el perfil TUBO 70$\times$70$\times$4 mm
          presenta un ratio de utilizacion maximo de \textbf{0.67} — muy por debajo
          del limite de 1.0 segun AISC 360-10. \textbf{Todos los miembros OK.}

    \item El desplazamiento maximo ($\approx 43$ mm, $L/385$) es aceptable
          para una cubierta sin requisitos estrictos de servicio.

    \item Las conexiones se modelaron rigidas. En la practica, las uniones
          soldadas de la cercha se aproximan a este comportamiento.

    \item Verificar uniones soldadas y placa base segun NSR-10 Capitulo F.2.
\end{enumerate}
}

\vspace{1.5cm}
\begin{center}
\begin{tikzpicture}
  \fill[primary] (0,0) rectangle (\textwidth,0.8pt);
  \fill[accent] (0.35\textwidth,0) rectangle (0.65\textwidth,0.8pt);
\end{tikzpicture}
\end{center}

%============================================
\appendix
\section{Archivo de entrada de datos}
%============================================

{
\sffamily\small
Archivo de entrada en formato STAAD (.std) utilizado para el modelo.
Ancho aferente: 0.65 m.
}

\vspace{0.3cm}
{\scriptsize\ttfamily
\verbatiminput{../examples/cercha_warren.std}
}

\end{document}
